\section{Jets}
\label{sec:ch6_jets}

\subsection{Reconstruction}

Jets are reconstructed from the collimated hadrons produced through parton showering and hadronisation.
Particle-flow constituents are used to combine the charged-particle momentum measurements from the ID with calorimeter measurements of the remaining energy~\cite{ATLAS2017ParticleFlow}.
The main steps are:
\begin{enumerate}
    \item Well-measured tracks are matched to calorimeter topological clusters.
    \item The expected calorimeter energy deposited by the matched charged particles is subtracted from the clusters, avoiding double counting when the track momenta are used.
    \item The selected tracks and remaining calorimeter contributions are retained as the jet inputs. Vertex association is used to suppress charged-particle contributions from pile-up.
\end{enumerate}
Calorimeter deposits from neutral particles and charged particles without a selected track remain in the calorimeter contribution.
The input clusters are initially measured at the electromagnetic energy scale, and the hadronic response is corrected by the jet calibration.

The constituents are clustered with the anti-$k_t$ algorithm using $R=0.4$~\cite{Cacciari2008AntiKt}.
At each iteration, pairwise and beam distances are evaluated:
\begin{align}
    d_{ij}&=\min\!\left(p_{\mathrm{T},i}^{-2},p_{\mathrm{T},j}^{-2}\right)
       \frac{\Delta R_{y,ij}^{\,2}}{R^2},
    &d_{iB}&=p_{\mathrm{T},i}^{-2},
    \label{eq:ch6_antikt}\\
    \Delta R_{y,ij}&=\sqrt{(y_i-y_j)^2+(\phi_i-\phi_j)^2},
    &y&=\frac{1}{2}\ln\frac{E+p_z}{E-p_z}.
    \label{eq:ch6_rapidity}
\end{align}
Here, $y$ is rapidity and $B$ labels the beam distance.
If the smallest distance is $d_{ij}$, the two constituent four-momenta are combined.
If it is $d_{iB}$, constituent $i$ is declared a jet and removed from the sequence.
The procedure is repeated until no constituents remain.
The inverse-$\pT^2$ weighting favours the clustering of soft radiation around hard constituents, giving isolated jets approximately circular boundaries.
The algorithm is infrared and collinear safe: the hard-jet configuration is unchanged by an arbitrarily soft emission or a collinear splitting.
The rapidity distance $\Delta R_y$ differs from the pseudorapidity distance $\Delta R$ defined in Eq.~\eqref{eq:delta_r}.
The two distances coincide in the massless limit.

\subsection{Identification and pile-up rejection}

Two quality requirements are applied:
\begin{itemize}
    \item \textbf{Jet cleaning.} The \texttt{LooseBad} criteria are used to reject jets with energy deposits characteristic of calorimeter noise or non-collision backgrounds.
    \item \textbf{Pile-up rejection.} The neural-network jet-vertex tagger, NNJVT, is used with the \texttt{FixedEffPt} working point. Track-to-vertex association is used to distinguish jets from the hard scattering from jets produced in pile-up collisions.
\end{itemize}
The NNJVT requirement is applied within the kinematic range of its calibration.

\subsection{Calibration}

The reconstructed jet four-momentum is calibrated through the following sequence~\cite{ATLAS2021JetCalibration}:
\begin{enumerate}
    \item \textbf{Pile-up subtraction.} The average additional transverse momentum is subtracted using the event transverse-momentum density $\rho$ and jet area $A$, giving a correction proportional to $\rho A$. Residual dependence on the pile-up activity is then corrected.
    \item \textbf{Simulation-based energy and direction calibration.} Reconstructed jets are compared with particle-level jets in simulation to correct the average energy response and pseudorapidity bias.
    \item \textbf{Global sequential calibration.} Dependence on the particle composition and shower development is reduced using tracking and calorimeter observables sensitive to variations in the jet response.
    \item \textbf{Residual in situ calibration.} Remaining differences between data and simulation are measured with momentum balance in $Z$+jet, photon+jet and multijet events. This final correction is applied to collision data.
\end{enumerate}
Particle-level jets are built from simulated stable particles, excluding muons and neutrinos, with the corresponding jet algorithm and provide the reference for the response calibration.
Jet energy-scale and resolution uncertainties are propagated to the analysis selections and to $\met$.
Scale factors for the pile-up rejection efficiency are applied to simulated event weights.
